\section{Volumetric experimental methods and supporting results}
\label{sec:experiments}

This appendix supplies the methods and the additional results behind the evaluation of the first application in Section~\ref{main:evidence}. Subsections~\ref{sec:implementation} and~\ref{sec:experiments-setup} give the implementation, metrics, and baselines. Subsection~\ref{sec:experiments-geometries} reports all fifteen whole-part runs. Subsection~\ref{sec:experiments-resolution} separates the two non-convex failure mechanisms. Subsection~\ref{sec:experiments-pieces-independent} documents the partitions, hull-gap evaluation, and same-budget comparison. The methods of the kinematics-first experiment are in Appendix~\ref{app:kin}.

\subsection{Implementation}
\label{sec:implementation}

\paragraph{Geometry.}
A watertight polygonal solid is scaled into $[0,1]^3$ with a margin of $0.05$, and all distances are in these normalised units. The solid is divided into tetrahedra by TetGen \cite{si2015tetgen} without changing its surface. The sites are the $n=8^m$ cell centres of the order-$m$ Hilbert curve, generated with Skilling's integer transform \cite{skilling2004}; for orders 1 to 6 we checked the face-adjacent, nested ordering that Subsection~\ref{sec:vsfc-hilbert} assumes. No waypoint is projected, removed, or smoothed.

\paragraph{Solver.}
The equal-mass weights \eqref{eq:main-equal-mass} are computed by the unmodified semi-discrete solver of Geogram~1.10.1 \cite{levy2015numerical}. The solve runs coarse to fine: at order two, initial weights are chosen so that every cell has positive mass, and Newton iterations \cite{kitagawa2019convergence} follow; each higher order starts from the previous order's weights, interpolated onto the nested grid. The stopping threshold bounds each relative cell-mass error by $10^{-6}$. The centroids \eqref{eq:main-centroid} are integrated exactly over the clipped cells with a degree-two tetrahedral rule.

\paragraph{Acceptance.}
A result is accepted only if every relative cell-mass error is below $10^{-6}$, no cell is empty, every centroid is finite, and every centroid is power-nearest to its own site, which confirms that the waypoints keep their Hilbert labels (Proposition~\ref{prop:sd-mass}(ii)). Every transport solve reported here passed; acceptance does not certify geometric containment. On the unit cube, the solver returned zero weights and centroids equal to the sites to $10^{-15}$, as Proposition~\ref{prop:sd-mass}(v) requires.

\paragraph{Evaluation domain.}
In the whole-part series, containment and segment metrics for the cylinder, sphere, and torus use the analytic solids; the decomposition experiments use the polygonal meshes.

\subsection{Metrics and baselines}
\label{sec:experiments-setup}

\emph{Cloud metrics} ignore the order of the waypoints. Containment is the fraction of waypoints inside $\Omega$. Coverage is the fraction of voxels of a $20^3$ grid that meet $\Omega$ and contain a waypoint, and its coefficient of variation (CV) is that of the waypoint counts over those voxels, lower being more uniform. The cell diameter $D_m$ is the largest diameter of a clipped Laguerre cell. \emph{Sequence metrics} depend on the order. Jumps count steps longer than five times the mean step; this is a resolution-dependent indicator, not a continuity test. \emph{Segment metrics} sample each segment at a spacing of at most $10^{-3}$; the excursion depth is the distance from an exterior sample to the surface. The covering radius is the largest distance from a point of the evaluation sample to its nearest waypoint.

The parts are a cylinder, a sphere, a solid torus, an L-bracket with a reentrant edge, and a complex part consisting of a block with several holes and a boss (genus three). The evaluation sample is a fixed, seeded set of $10\,000$ uniform interior points; the map itself involves no random choice. An independent uniform sample of the same size provides a coverage reference.

The layered baseline visits the cells of the grid of side $2^{-m}$ whose centres lie in $\Omega$, layer by layer, traversing each layer row by row in alternating directions; it is the planar slicer with $h=2^{-m}$. Every move between cells that do not share a face is a travel move, hence a stop, and moves between layers are counted separately. For the transported path, each maximal \emph{exterior run} of segments leaving $\Omega$ counts as one stop.

\subsection{Whole-part runs}
\label{sec:experiments-geometries}

\paragraph{Setup.}
Each of the five parts is solved as a single target with uniform density at orders four, five, and six, giving fifteen runs.

\paragraph{Containment and coverage.}
Table~\ref{tab:experiments-geometries} reports every uniform run. On the torus at order four, coverage is $0.882$ with CV $0.613$, against $0.861$ and $0.732$ for the uniform reference. Re-evaluating on ten independent samples changes coverage, CV, and covering radius by at most $0.006$. Convex containment follows from Proposition~\ref{prop:sd-mass}(iv); jump counts and step lengths are measured.

% Written by python -m mapchain tables from results/evaluation/whole_part.
\begin{papertable}
  \begin{papertabular}{lrrrrrrrrrr}
    \toprule
    Part & $m$ & Inside & \papertablehead{Ext.\\seg.} & \papertablehead{Ext.\\length} & Depth & \papertablehead{Ext.\\runs} & Jumps & $D_m$ & \papertablehead{Cover.\\radius} & \papertablehead{Max.\\step} \\
    \midrule
    Cylinder & 4 & $1.000$ & $0$ & $0$ & $0$ & $0$ & $0$ & $0.120$ & $0.047$ & $0.069$ \\
    & 5 & $1.000$ & $0$ & $0$ & $0$ & $0$ & $0$ & $0.065$ & $0.023$ & $0.039$ \\
    & 6 & $1.000$ & $0$ & $0$ & $0$ & $0$ & $0$ & $0.035$ & $0.012$ & $0.023$ \\
    Sphere & 4 & $1.000$ & $0$ & $0$ & $0$ & $0$ & $0$ & $0.120$ & $0.043$ & $0.073$ \\
    & 5 & $1.000$ & $0$ & $0$ & $0$ & $0$ & $0$ & $0.066$ & $0.022$ & $0.043$ \\
    & 6 & $1.000$ & $0$ & $0$ & $0$ & $0$ & $0$ & $0.036$ & $0.012$ & $0.025$ \\
    Torus & 4 & $1.000$ & $<0.001$ & $0.001$ & $0.029$ & $1$ & $0$ & $0.179$ & $0.039$ & $0.205$ \\
    & 5 & $1.000$ & $0.003$ & $0.004$ & $0.032$ & $53$ & $1$ & $0.105$ & $0.024$ & $0.198$ \\
    & 6 & $0.999$ & $0.002$ & $0.008$ & $0.032$ & $241$ & $255$ & $0.092$ & $0.014$ & $0.195$ \\
    L-bracket & 4 & $0.973$ & $0.049$ & $0.182$ & $0.264$ & $53$ & $97$ & $0.858$ & $0.062$ & $0.396$ \\
    & 5 & $0.985$ & $0.026$ & $0.169$ & $0.280$ & $224$ & $477$ & $0.853$ & $0.033$ & $0.405$ \\
    & 6 & $0.993$ & $0.012$ & $0.157$ & $0.289$ & $836$ & $2201$ & $0.850$ & $0.017$ & $0.412$ \\
    Complex part & 4 & $0.988$ & $0.038$ & $0.058$ & $0.045$ & $95$ & $0$ & $0.197$ & $0.041$ & $0.166$ \\
    & 5 & $0.998$ & $0.016$ & $0.053$ & $0.089$ & $323$ & $174$ & $0.167$ & $0.022$ & $0.182$ \\
    & 6 & $0.999$ & $0.007$ & $0.033$ & $0.083$ & $1103$ & $761$ & $0.165$ & $0.012$ & $0.181$ \\
    \bottomrule
  \end{papertabular}
  \caption{All fifteen uniform runs, at $8^m$ waypoints (4096, $32\,768$, and $262\,144$). ``Inside'' is the fraction of waypoints inside $\Omega$; ``Ext.\ seg.'' and ``Ext.\ length'' are the fractions of exterior segments and exterior path length; ``Depth'' is the deepest sampled excursion. Distances ($D_m$, depth, covering radius, and maximum step) are in normalised units (Appendix~\ref{sec:implementation}); exterior runs and jumps are counts.}
  \label{tab:experiments-geometries}
\end{papertable}


\paragraph{Stops against the layered baseline.}
\label{sec:experiments-interruptions}
Table~\ref{tab:main-stops} gives the counts. The L-bracket's cross-sections are rectangles aligned with the raster grid, so its rows almost always continue within a layer.

\subsection{Non-convex parts and refinement}
\label{sec:experiments-resolution}
\label{sec:sd-nonconvex}

These results reuse the whole-part runs and support the main-text claim that refinement reduces but does not remove the non-convex failures.

\paragraph{Outside centroids and crossing chords.}
The two non-convex mechanisms of Section~\ref{main:semidiscrete} can be told apart by asking whether an exterior segment has both endpoints inside. On the L-bracket, none does at any order: every excursion starts or ends at a centroid in the region between the arms, where cells clipped into two arms place their centroids. On the complex part, a growing share have both endpoints inside (from $65$ of $156$ exterior segments at order four to $1218$ of $1765$ at order six): consecutive centroids lie on opposite sides of a hole and the chord between them crosses it. The torus lies between the two.

\subsection{Decomposed source}
\label{sec:experiments-pieces-independent}

These runs test Proposition~\ref{prop:sd-hull-gap} and Corollary~\ref{cor:sd-piece-hull-gap} on the L-bracket's two-box control and on an automatic partition of the polygonal torus. No path segments are added between pieces.

\paragraph{Setup: partitions and solves.}
TetGen, through the \texttt{tetgen}~0.8.4 Python package, meshes each object once (switches \texttt{pq1.4Q}). The L-bracket is split exactly into its two boxes. For the torus, CoACD~1.0.14 \cite{wei2022coacd} supplies convex hulls at threshold $0.05$ and seed zero. Each tetrahedron is assigned to the hull it lies deepest inside, connected components become pieces, and components below $1\%$ of total volume merge into a neighbour: this gives eight torus pieces. The resulting pieces need not be convex. Every tetrahedron has exactly one owner, and piece volumes conserve the total to rounding error.

All targets are uniform. Each piece receives order four, 4096 Hilbert sites, its conditional density, and time share $w_j$ equal to its volume fraction. The unchanged pipeline and acceptance checks of Subsection~\ref{sec:implementation} are used, and all ten solves pass.

\paragraph{Setup: hull-gap bounds.}
The two boxes have $\delta_j=0$. For each automatic piece, a branch-and-bound computation covers its convex hull by cubes, bounding the distance to the piece on each cube and refining until the upper and lower bounds $U_j$ and $L_j$ are within $2\times10^{-3}$. The computation uses floating-point geometry and winding numbers for inside tests, without an interval-arithmetic proof, so the bounds are numerical rather than exact. Their positive lower bounds show that the automatic pieces are not convex.

\paragraph{Results against the assigned pieces.}
Each segment is sampled at a spacing of at most $10^{-3}$ and counted exterior if any sample lies outside its piece; depth is the largest sampled distance. Table~\ref{tab:decomposition} reports the two cases; every sampled excursion is below its own $U_j$.

% Written by python -m mapchain tables from results/evaluation/pieces/summary.json.
\begin{papertable}
  \begin{papertabular}{lrrrrrr}
    \toprule
    Object & Pieces & Waypoints & $\max_j\gamma^{\mathrm{ub}}(\Omega_j)$ & In piece & \papertablehead{Exterior\\segments} & \papertablehead{Sampled\\depth} \\
    \midrule
    L-bracket & 2 & 8192 & $0$ & 100.00\% & 0.00\% & $0.0000$ \\
    Torus & 8 & $32\,768$ & $0.0563$ & 98.40\% & 5.21\% & $0.0232$ \\
    \bottomrule
  \end{papertabular}
  \caption{Independent piece paths, uniform density, 4096 waypoints per piece. ``In piece'' is the percentage of waypoints inside their assigned piece; exterior segments are measured against that same piece. Depth is the largest sampled excursion; $\gamma^{\mathrm{ub}}(\Omega_j)$ is the numerical upper bound on the hull gap of piece $j$. Distances are in normalised units (Appendix~\ref{sec:implementation}).}
  \label{tab:decomposition}
\end{papertable}


\paragraph{Results against the whole object.}
\label{sec:experiments-torus-comparison}
A same-budget comparison re-evaluates two saved torus runs with $32\,768$ waypoints each, the whole-part order-five run and the eight-piece order-four run, against the polygonal torus. Coverage, CV, and covering radius use the same $10\,000$ interior samples for both; each piece waypoint is weighted by its piece's time share, $w_j/4096$, since ignoring these weights would give a misleading CV of $0.84$ rather than $0.40$.

Table~\ref{tab:torus-whole-vs-pieces} reports measurements against both the assigned pieces and the whole torus; these containment tests use different domains.

\paragraph{Figure run at the whole-part budget.}
\label{sec:experiments-torus-figure}
The torus piece paths in Figure~\ref{fig:containment-repair} come from a separate run with the same partition, pipeline, and acceptance checks at order three: eight pieces of 512 waypoints, 4096 in total, the budget of the order-four whole-part path beside them. All eight solves pass. Against the assigned pieces, 334 of the 4088 segments ($8.2\%$) leave their own piece, with a deepest sampled excursion of $0.019$, below every $U_j$. Against the analytic torus, with 60 samples per segment, no segment leaves the part, whereas the order-four whole-part path has one exterior segment.

\paragraph{Limitations.}
Zero detected exterior segments is a sampled result, not a containment certificate, and the hull-gap bounds are numerical. The effect of assigned-piece crossings on local density and overlap is unmeasured, and the comparison concerns these two torus runs only.


\paragraph{Teaser provenance.}
Figure~\ref{fig:teaser} uses the Stanford bunny redistributed with libigl's tutorial data (6966 triangles), without repair or simplification. TetGen (\texttt{tetgen}~0.8.4, \texttt{pq1.4Q}) meshes it once; CoACD~1.0.14 at threshold $0.1$ and seed zero, followed by the tetrahedron assignment and $1\%$ component-merge rule above, gives nine pieces. Each uniform piece receives 4096 sites, its conditional density, and a time share equal to its volume fraction, so each waypoint has global mass $w_j/4096$. All nine piece solves and the additional whole-part order-five solve pass the acceptance checks. The upper teaser row uses the whole-part path and the lower row the piece paths. Records of transforms, parameters, versions, and checks are in \path{results/volumetric/decomposition/}.
